%%%%%%%%%%%%%%%%%%%%%%%%%%%%%%%%%%%%%%%%%%%%%%%%%%%%%%%%%%%%%%%%%%%%%%%%%%%%%%%%
\chapter{Alternative Stochastic Approaches to Quantum Theory}
\label{ch:AlternativeStochastic}
%%%%%%%%%%%%%%%%%%%%%%%%%%%%%%%%%%%%%%%%%%%%%%%%%%%%%%%%%%%%%%%%%%%%%%%%%%%%%%%%

%%%%%%%%%%%%%%%%%%%%%%%%%%%%%%%%%%%%%%%%%%%%%%%%%%%%%%%%%%%%%%%%%%%%%%%%%%%%%%%%
\section{Introduction}
%%%%%%%%%%%%%%%%%%%%%%%%%%%%%%%%%%%%%%%%%%%%%%%%%%%%%%%%%%%%%%%%%%%%%%%%%%%%%%%%

The remarkable predictive success of quantum mechanics has never
eliminated interest in the possibility that it may emerge from a deeper,
more fundamental theory. Since the earliest days of quantum mechanics,
many physicists have questioned whether the probabilistic character of
the theory is truly fundamental or merely reflects our ignorance of
underlying microscopic dynamics.

One possible answer is provided by hidden-variable theories, of which
the de Broglie--Bohm theory is the best-known example. Another,
conceptually different, possibility is that quantum behavior arises from
an underlying stochastic process. In this picture the apparent
indeterminacy of quantum mechanics does not originate from intrinsic
probability, but from continual interaction with microscopic random
degrees of freedom.

The idea has appeared repeatedly during the history of physics in
several independent forms. Although these approaches differ
substantially in their mathematical formulation, they all attempt to
answer essentially the same question:

\begin{quote}
\emph{Can quantum mechanics be understood as the statistical description
of an underlying stochastic dynamics?}
\end{quote}

This question has motivated a surprisingly diverse collection of
theories.

Some authors assume that particles themselves execute Brownian motion.
Others attribute the stochasticity to classical electromagnetic fields,
to random space-time geometry, to non-differentiable trajectories, to
matrix-valued microscopic dynamics, or even to stochastic evolution in
an auxiliary mathematical time.

Consequently, the term \emph{stochastic quantum theory} does not denote
a single well-defined framework but rather an entire family of
approaches sharing a common philosophical motivation while differing
greatly in physical assumptions.

The purpose of the present chapter is not merely to review these
theories individually. Equally important is to compare them
systematically, identify their common conceptual structure, and clarify
the physical questions they attempt to answer.

%%%%%%%%%%%%%%%%%%%%%%%%%%%%%%%%%%%%%%%%%%%%%%%%%%%%%%%%%%%%%%%%%%%%%%%%%%%%%%%%
\subsection{What Can Be Stochastic?}
%%%%%%%%%%%%%%%%%%%%%%%%%%%%%%%%%%%%%%%%%%%%%%%%%%%%%%%%%%%%%%%%%%%%%%%%%%%%%%%%

One of the principal sources of confusion in the literature is that the
word \emph{stochastic} is used in several fundamentally different ways.

The first possibility is that the particle trajectory itself is random.
This viewpoint is adopted in Nelson's stochastic mechanics.

A second possibility is that the particle remains completely classical,
while the electromagnetic field possesses a universal stochastic
background. This is the central hypothesis of stochastic
electrodynamics.

A third possibility is that the stochastic process takes place not in
physical space-time but in the infinite-dimensional space of field
configurations. This idea forms the basis of Parisi--Wu stochastic
quantization.

Other approaches attribute randomness to the geometry of space-time
itself, to fractal properties of trajectories, or to the statistical
mechanics of noncommuting matrices.

Thus, before comparing the various theories, one must first answer a
more fundamental question:

\begin{center}
\emph{What physical quantity is assumed to fluctuate?}
\end{center}

As we shall see, the answer to this question largely determines the
mathematical structure and physical interpretation of each theory.

%%%%%%%%%%%%%%%%%%%%%%%%%%%%%%%%%%%%%%%%%%%%%%%%%%%%%%%%%%%%%%%%%%%%%%%%%%%%%%%%
\subsection{Common Objectives}
%%%%%%%%%%%%%%%%%%%%%%%%%%%%%%%%%%%%%%%%%%%%%%%%%%%%%%%%%%%%%%%%%%%%%%%%%%%%%%%%

Despite their diversity, nearly all stochastic approaches pursue several
closely related goals.

These include

\begin{itemize}

\item explaining the origin of Planck's constant,

\item deriving the Schrödinger equation from more primitive dynamical
principles,

\item understanding the physical meaning of the wave function,

\item explaining quantum probabilities,

\item clarifying the measurement problem,

\item establishing a deeper connection between classical and quantum
physics.

\end{itemize}

Not every theory attempts to solve all of these problems.

For example, stochastic quantization is primarily a computational method
for constructing quantum field theories rather than an interpretation of
quantum mechanics.

Conversely, stochastic electrodynamics explicitly seeks a classical
physical origin of quantum phenomena.

Recognizing these different objectives is essential for a fair
comparison.

%%%%%%%%%%%%%%%%%%%%%%%%%%%%%%%%%%%%%%%%%%%%%%%%%%%%%%%%%%%%%%%%%%%%%%%%%%%%%%%%
\subsection{Criteria for Comparison}
%%%%%%%%%%%%%%%%%%%%%%%%%%%%%%%%%%%%%%%%%%%%%%%%%%%%%%%%%%%%%%%%%%%%%%%%%%%%%%%%

Throughout this chapter we shall compare the various theories according
to a common set of questions.

For each approach we shall ask:

\begin{enumerate}

\item What is the fundamental stochastic variable?

\item What is the evolution parameter?

\item Is the theory relativistic?

\item Does the theory explain the origin of Planck's constant?

\item Can it derive the Schrödinger equation?

\item Can it describe spin?

\item Does it naturally generalize to quantum field theory?

\item How are gauge interactions incorporated?

\item Is gravity included?

\item What are its principal successes?

\item What are its most serious unresolved problems?

\end{enumerate}

Using identical criteria for every theory allows a meaningful comparison
that is largely absent from the existing literature.

%%%%%%%%%%%%%%%%%%%%%%%%%%%%%%%%%%%%%%%%%%%%%%%%%%%%%%%%%%%%%%%%%%%%%%%%%%%%%%%%
\subsection{Organization of This Chapter}
%%%%%%%%%%%%%%%%%%%%%%%%%%%%%%%%%%%%%%%%%%%%%%%%%%%%%%%%%%%%%%%%%%%%%%%%%%%%%%%%

The chapter begins with stochastic electrodynamics, the oldest and most
extensively developed classical stochastic approach to quantum theory.
We then review stochastic quantization, scale relativity, trace
dynamics, and several more recent proposals.

The chapter concludes with a systematic comparison of all approaches and
a discussion of the lessons they provide for the stochastic
Hamiltonian framework developed in the remainder of this monograph.

%%%%%%%%%%%%%%%%%%%%%%%%%%%%%%%%%%%%%%%%%%%%%%%%%%%%%%%%%%%%%%%%%%%%%%%%%%%%%%%%
\subsection*{Editorial Comment}
%%%%%%%%%%%%%%%%%%%%%%%%%%%%%%%%%%%%%%%%%%%%%%%%%%%%%%%%%%%%%%%%%%%%%%%%%%%%%%%%

The diversity of stochastic approaches reviewed in this chapter should
not be regarded as evidence of failure. Rather, it reflects the
difficulty of the underlying problem.

Quantum mechanics is extraordinarily successful experimentally, and any
attempt to derive it from deeper principles faces exceptionally
stringent constraints. The existence of multiple stochastic programs,
developed largely independently over more than half a century,
demonstrates that the possibility of an underlying stochastic dynamics
has remained scientifically compelling despite the absence of a
universally accepted theory.

Our objective is therefore not to advocate a particular approach, but to
understand what each contributes, where each succeeds, where each fails,
and which ideas appear sufficiently robust to serve as ingredients of a
more complete relativistic stochastic theory.



%%%%%%%%%%%%%%%%%%%%%%%%%%%%%%%%%%%%%%%%%%%%%%%%%%%%%%%%%%%%%%%%%%%%%%%%%%%%%%%%
\section{Historical Development of Stochastic Electrodynamics}
\label{sec:SEDHistory}
%%%%%%%%%%%%%%%%%%%%%%%%%%%%%%%%%%%%%%%%%%%%%%%%%%%%%%%%%%%%%%%%%%%%%%%%%%%%%%%%

Stochastic electrodynamics (SED) is one of the oldest attempts to
understand quantum phenomena in terms of an underlying classical
stochastic dynamics. Unlike stochastic mechanics, where randomness is
attributed to the motion of particles, SED assumes that the
electromagnetic field itself possesses an intrinsic random component
which exists even at absolute zero temperature.

From this viewpoint particles always obey classical equations of motion.
Their apparently quantum behavior results from continual interaction
with an omnipresent stochastic electromagnetic background.

The central hypothesis may therefore be summarized in a single sentence:

\begin{quote}
Quantum phenomena originate not from quantization of matter, but from
the interaction of classical matter with a universal classical
stochastic radiation field.
\end{quote}

This idea has a long and fascinating history extending over more than a
century.

%%%%%%%%%%%%%%%%%%%%%%%%%%%%%%%%%%%%%%%%%%%%%%%%%%%%%%%%%%%%%%%%%%%%%%%%%%%%%%%%
\subsection{Origins in Classical Radiation Theory}
%%%%%%%%%%%%%%%%%%%%%%%%%%%%%%%%%%%%%%%%%%%%%%%%%%%%%%%%%%%%%%%%%%%%%%%%%%%%%%%%

The roots of SED lie in the crisis of classical radiation theory at the
end of the nineteenth century.

Maxwell's equations had achieved spectacular success in describing
electromagnetic phenomena. At the same time, classical statistical
mechanics appeared capable of explaining the thermodynamic properties of
matter.

The difficulty arose when these two theories were combined to describe
thermal radiation.

According to classical statistical mechanics, every normal mode of the
electromagnetic field should possess the average energy

\[
\langle E\rangle=k_BT.
\]

Since the number of modes increases rapidly with frequency, the total
energy diverges,

\[
U=\int_0^\infty u(\omega)d\omega=\infty,
\]

producing the well-known ultraviolet catastrophe.

Planck resolved this difficulty in 1900 by introducing energy quanta,

\[
E=n\hbar\omega.
\]

Although enormously successful, this proposal raised a profound
question.

Is quantization truly fundamental, or does it merely provide an
effective description of some deeper classical dynamics?

This question eventually became one of the principal motivations for
stochastic electrodynamics.

%%%%%%%%%%%%%%%%%%%%%%%%%%%%%%%%%%%%%%%%%%%%%%%%%%%%%%%%%%%%%%%%%%%%%%%%%%%%%%%%
\subsection{Einstein, Lorentz and Classical Fluctuations}
%%%%%%%%%%%%%%%%%%%%%%%%%%%%%%%%%%%%%%%%%%%%%%%%%%%%%%%%%%%%%%%%%%%%%%%%%%%%%%%%

Long before the modern formulation of SED, several pioneers had already
recognized the importance of fluctuations in electromagnetic theory.

Einstein's work on Brownian motion demonstrated that random microscopic
forces could produce observable macroscopic phenomena.

Lorentz investigated fluctuations generated by classical
electromagnetic fields.

Planck himself repeatedly reconsidered the possibility that the vacuum
might possess an intrinsic electromagnetic energy even at zero
temperature.

Although none of these ideas constituted stochastic electrodynamics in
its modern form, they established the conceptual foundations upon which
the later theory was built.

%%%%%%%%%%%%%%%%%%%%%%%%%%%%%%%%%%%%%%%%%%%%%%%%%%%%%%%%%%%%%%%%%%%%%%%%%%%%%%%%
\subsection{The Modern Beginning}
%%%%%%%%%%%%%%%%%%%%%%%%%%%%%%%%%%%%%%%%%%%%%%%%%%%%%%%%%%%%%%%%%%%%%%%%%%%%%%%%

Modern stochastic electrodynamics emerged during the 1960s.

The decisive step was the realization that Maxwell's equations admit a
Lorentz-invariant random solution whose spectral density is proportional
to

\[
\omega^3.
\]

Instead of quantizing the electromagnetic field, one assumes that this
random radiation field exists as a classical solution of Maxwell's
equations.

Among the principal contributors were

\begin{itemize}

\item Timothy H. Boyer,

\item Trevor W. Marshall,

\item Luis de la Peña,

\item Ana María Cetto,

\item Francisco Francia,

\item Theodore Puthoff,

\item Daniel Cole,

\item Yorick Savoy,

\item Gerhard Hegerfeldt (critical analyses),

\item and many others.

\end{itemize}

These authors transformed SED from an isolated idea into a coherent
research program.

%%%%%%%%%%%%%%%%%%%%%%%%%%%%%%%%%%%%%%%%%%%%%%%%%%%%%%%%%%%%%%%%%%%%%%%%%%%%%%%%
\subsection{The Boyer Revolution}
%%%%%%%%%%%%%%%%%%%%%%%%%%%%%%%%%%%%%%%%%%%%%%%%%%%%%%%%%%%%%%%%%%%%%%%%%%%%%%%%

The work of Boyer occupies a special position in the history of SED.

He demonstrated that many phenomena usually regarded as intrinsically
quantum could be reproduced by combining classical electrodynamics with
a Lorentz-invariant stochastic zero-point radiation field.

Among his best-known achievements are

\begin{itemize}

\item derivation of the Planck spectrum,

\item blackbody radiation,

\item harmonic oscillator,

\item Casimir forces,

\item diamagnetism,

\item thermal equilibrium,

\item Unruh-type phenomena.

\end{itemize}

Whether these results constitute a genuine derivation of quantum
mechanics remains controversial.

Nevertheless, Boyer's work demonstrated that classical physics is far
more powerful than previously believed once stochastic vacuum radiation
is included.

%%%%%%%%%%%%%%%%%%%%%%%%%%%%%%%%%%%%%%%%%%%%%%%%%%%%%%%%%%%%%%%%%%%%%%%%%%%%%%%%
\subsection{The Mexican School}
%%%%%%%%%%%%%%%%%%%%%%%%%%%%%%%%%%%%%%%%%%%%%%%%%%%%%%%%%%%%%%%%%%%%%%%%%%%%%%%%

A second major center of SED developed around Luis de la Peña and Ana
María Cetto.

Their approach emphasized the statistical mechanics of particles
interacting continuously with the zero-point field.

Particular attention was devoted to

\begin{itemize}

\item derivation of the Schrödinger equation,

\item interpretation of the wave function,

\item momentum fluctuations,

\item diffusion processes,

\item relation to Nelson's mechanics,

\item foundations of quantum mechanics.

\end{itemize}

Their comprehensive monograph remains one of the principal references in
the field.

%%%%%%%%%%%%%%%%%%%%%%%%%%%%%%%%%%%%%%%%%%%%%%%%%%%%%%%%%%%%%%%%%%%%%%%%%%%%%%%%
\subsection{Revival During the Twenty-First Century}
%%%%%%%%%%%%%%%%%%%%%%%%%%%%%%%%%%%%%%%%%%%%%%%%%%%%%%%%%%%%%%%%%%%%%%%%%%%%%%%%

Interest in SED increased again after approximately 2000 for several
reasons.

First,

high-precision numerical simulations became possible.

Second,

renewed interest in quantum foundations encouraged reconsideration of
alternative formulations of quantum mechanics.

Third,

connections emerged between SED and several modern topics including

\begin{itemize}

\item Casimir physics,

\item vacuum fluctuations,

\item nonequilibrium statistical mechanics,

\item quantum thermodynamics,

\item relativistic stochastic dynamics.

\end{itemize}

Although SED remains outside the mainstream interpretation of quantum
mechanics, it continues to generate active research.

%%%%%%%%%%%%%%%%%%%%%%%%%%%%%%%%%%%%%%%%%%%%%%%%%%%%%%%%%%%%%%%%%%%%%%%%%%%%%%%%
\subsection{Relation to Other Stochastic Theories}
%%%%%%%%%%%%%%%%%%%%%%%%%%%%%%%%%%%%%%%%%%%%%%%%%%%%%%%%%%%%%%%%%%%%%%%%%%%%%%%%

Historically, SED developed almost independently of Nelson's stochastic
mechanics.

This independence is rather surprising because both theories seek a
classical stochastic origin of quantum behavior.

The principal distinction lies in the source of stochasticity.

Nelson assumes

\[
\boxed{
\text{Particle motion}
\longrightarrow
\text{stochastic}.
}
\]

SED assumes

\[
\boxed{
\text{Electromagnetic field}
\longrightarrow
\text{stochastic}.
}
\]

This apparently simple difference has profound consequences for the
mathematical structure of the resulting theories.

%%%%%%%%%%%%%%%%%%%%%%%%%%%%%%%%%%%%%%%%%%%%%%%%%%%%%%%%%%%%%%%%%%%%%%%%%%%%%%%%
\subsection*{Critical Assessment}
%%%%%%%%%%%%%%%%%%%%%%%%%%%%%%%%%%%%%%%%%%%%%%%%%%%%%%%%%%%%%%%%%%%%%%%%%%%%%%%%

The historical development of stochastic electrodynamics illustrates an
important lesson.

SED should not be viewed merely as an unsuccessful attempt to replace
quantum mechanics.

Rather, it represents a sustained effort to determine how much of
quantum physics can be recovered from classical electrodynamics when the
vacuum is regarded as a genuine stochastic physical medium.

Whether or not this program ultimately succeeds is less important than
the insight it provides into the relationship between classical field
theory, stochastic processes and quantum phenomena.

From the viewpoint of the present monograph, SED constitutes one of the
principal alternative stochastic programs that any new relativistic
stochastic theory must carefully compare itself with.


%%%%%%%%%%%%%%%%%%%%%%%%%%%%%%%%%%%%%%%%%%%%%%%%%%%%%%%%%%%%%%%%%%%%%%%%%%%%%%%%
\section{Mathematical Foundations of Classical Zero-Point Radiation}
\label{sec:SEDMathematical}
%%%%%%%%%%%%%%%%%%%%%%%%%%%%%%%%%%%%%%%%%%%%%%%%%%%%%%%%%%%%%%%%%%%%%%%%%%%%%%%%

The defining postulate of stochastic electrodynamics is the existence of
a classical random electromagnetic field that persists even at absolute
zero temperature. Unlike ordinary thermal radiation, this field is
assumed to be a permanent property of the physical vacuum.

The particles themselves remain entirely classical. Their stochastic
motion arises solely through interaction with this universal
electromagnetic background.

Consequently, the mathematical formulation of SED begins not with
particle dynamics but with the statistical properties of solutions of
Maxwell's equations.

%%%%%%%%%%%%%%%%%%%%%%%%%%%%%%%%%%%%%%%%%%%%%%%%%%%%%%%%%%%%%%%%%%%%%%%%%%%%%%%%
\subsection{Free Electromagnetic Field}
%%%%%%%%%%%%%%%%%%%%%%%%%%%%%%%%%%%%%%%%%%%%%%%%%%%%%%%%%%%%%%%%%%%%%%%%%%%%%%%%

In the absence of charges and currents the Maxwell equations are

\begin{align}
\nabla\cdot\mathbf E &=0,
&
\nabla\cdot\mathbf B &=0,
\\
\nabla\times\mathbf E
&=
-\frac1c
\frac{\partial\mathbf B}{\partial t},
&
\nabla\times\mathbf B
&=
\frac1c
\frac{\partial\mathbf E}{\partial t}.
\end{align}

Introducing the four-potential

\[
A^\mu=(\phi,\mathbf A),
\]

and adopting the Lorenz gauge,

\[
\partial_\mu A^\mu=0,
\]

the field equations reduce to

\begin{equation}
\square A^\mu=0,
\label{eq:WaveEquation}
\end{equation}

where

\[
\square
=
\eta^{\mu\nu}\partial_\mu\partial_\nu
=
\frac1{c^2}
\frac{\partial^2}{\partial t^2}
-\nabla^2.
\]

Equation (\ref{eq:WaveEquation}) possesses infinitely many plane-wave
solutions.

SED assumes that the physical vacuum contains a random superposition of
all such modes.

%%%%%%%%%%%%%%%%%%%%%%%%%%%%%%%%%%%%%%%%%%%%%%%%%%%%%%%%%%%%%%%%%%%%%%%%%%%%%%%%
\subsection{Plane-Wave Expansion}
%%%%%%%%%%%%%%%%%%%%%%%%%%%%%%%%%%%%%%%%%%%%%%%%%%%%%%%%%%%%%%%%%%%%%%%%%%%%%%%%

The vector potential is expanded as

\begin{equation}
\mathbf A(\mathbf x,t)
=
\sum_{\lambda=1}^{2}
\int
\frac{d^3k}{(2\pi)^3}
\,
\boldsymbol{\epsilon}_{\lambda}(\mathbf k)
\,
a_{\lambda}(\mathbf k)
\cos
\left(
\mathbf k\cdot\mathbf x
-\omega t
+\phi_{\lambda}(\mathbf k)
\right),
\label{eq:RandomExpansion}
\end{equation}

where

\[
\omega=c|\mathbf k|,
\]

and

\[
\boldsymbol{\epsilon}_{\lambda}\cdot\mathbf k=0
\]

are transverse polarization vectors.

The amplitudes satisfy a prescribed spectral distribution, whereas the
phases

\[
\phi_{\lambda}(\mathbf k)
\]

are independent random variables uniformly distributed over

\[
0\le\phi<2\pi.
\]

This random-phase approximation is the central statistical assumption of
SED.

%%%%%%%%%%%%%%%%%%%%%%%%%%%%%%%%%%%%%%%%%%%%%%%%%%%%%%%%%%%%%%%%%%%%%%%%%%%%%%%%
\subsection{Random Variables}
%%%%%%%%%%%%%%%%%%%%%%%%%%%%%%%%%%%%%%%%%%%%%%%%%%%%%%%%%%%%%%%%%%%%%%%%%%%%%%%%

Ensemble averages are taken over the random phases.

Immediately,

\[
\langle
\cos(\theta+\phi)
\rangle
=
0,
\]

which implies

\begin{equation}
\boxed{
\langle\mathbf E\rangle
=
\langle\mathbf B\rangle
=
0.
}
\end{equation}

The vacuum therefore contains no preferred electric or magnetic field.

However,

quadratic averages do not vanish.

For example,

\[
\langle E_iE_j\rangle\neq0,
\]

indicating that the vacuum possesses finite electromagnetic
fluctuations despite zero mean fields.

This feature is completely analogous to ordinary Brownian motion, where

\[
\langle dW\rangle=0,
\qquad
\langle(dW)^2\rangle=dt.
\]

Thus,

the physically significant quantities are not the fields themselves but
their correlation functions.

%%%%%%%%%%%%%%%%%%%%%%%%%%%%%%%%%%%%%%%%%%%%%%%%%%%%%%%%%%%%%%%%%%%%%%%%%%%%%%%%
\subsection{Two-Point Correlation Functions}
%%%%%%%%%%%%%%%%%%%%%%%%%%%%%%%%%%%%%%%%%%%%%%%%%%%%%%%%%%%%%%%%%%%%%%%%%%%%%%%%

The statistical properties of the zero-point field are completely
specified by the two-point correlation functions,

\begin{equation}
C_{ij}(x,x')
=
\left<
E_i(x)
E_j(x')
\right>.
\end{equation}

For a homogeneous and isotropic vacuum,

the correlation function depends only upon the coordinate difference,

\[
x-x'.
\]

Translation invariance therefore implies

\[
C_{ij}(x,x')
=
C_{ij}(x-x').
\]

Similarly,

rotational invariance restricts the tensor structure to combinations of

\[
\delta_{ij},
\qquad
k_ik_j.
\]

Lorentz invariance places even stronger constraints and ultimately fixes
the spectral dependence of the vacuum field.

%%%%%%%%%%%%%%%%%%%%%%%%%%%%%%%%%%%%%%%%%%%%%%%%%%%%%%%%%%%%%%%%%%%%%%%%%%%%%%%%
\subsection{Energy Density}
%%%%%%%%%%%%%%%%%%%%%%%%%%%%%%%%%%%%%%%%%%%%%%%%%%%%%%%%%%%%%%%%%%%%%%%%%%%%%%%%

The electromagnetic energy density is

\begin{equation}
u
=
\frac1{8\pi}
(E^2+B^2)
\end{equation}

(in Gaussian units).

Taking the ensemble average gives

\[
\langle u\rangle
=
\frac1{8\pi}
\left<
E^2+B^2
\right>.
\]

Using the plane-wave expansion yields

\begin{equation}
\langle u\rangle
=
\int_0^\infty
u(\omega)
\,d\omega,
\end{equation}

where

\[
u(\omega)
\]

is the spectral energy density.

Determining this function constitutes the central problem of SED.

%%%%%%%%%%%%%%%%%%%%%%%%%%%%%%%%%%%%%%%%%%%%%%%%%%%%%%%%%%%%%%%%%%%%%%%%%%%%%%%%
\subsection{Why the Spectrum Is Not Arbitrary}
%%%%%%%%%%%%%%%%%%%%%%%%%%%%%%%%%%%%%%%%%%%%%%%%%%%%%%%%%%%%%%%%%%%%%%%%%%%%%%%%

At first sight one might imagine choosing any spectral density.

This is impossible.

The vacuum must satisfy simultaneously

\begin{enumerate}

\item Maxwell's equations,

\item statistical homogeneity,

\item isotropy,

\item Lorentz invariance,

\item scale invariance (at zero temperature),

\item positivity of the energy density.

\end{enumerate}

These requirements are so restrictive that only one spectral form
survives (up to an overall multiplicative constant).

As first emphasized by Boyer,

the Lorentz-invariant zero-temperature spectrum is proportional to

\[
\omega^3.
\]

The overall proportionality constant is then identified empirically with

\[
\frac{\hbar}{2\pi^2c^3},
\]

giving

\begin{equation}
\boxed{
u(\omega)
=
\frac{\hbar\omega^3}
{2\pi^2c^3}.
}
\label{eq:BoyerSpectrum}
\end{equation}

This result is among the most remarkable aspects of SED.

Planck's constant appears not because the field is quantized, but as the
measured intensity of a classical Lorentz-invariant stochastic vacuum.

%%%%%%%%%%%%%%%%%%%%%%%%%%%%%%%%%%%%%%%%%%%%%%%%%%%%%%%%%%%%%%%%%%%%%%%%%%%%%%%%
\subsection{The Origin of Planck's Constant}
%%%%%%%%%%%%%%%%%%%%%%%%%%%%%%%%%%%%%%%%%%%%%%%%%%%%%%%%%%%%%%%%%%%%%%%%%%%%%%%%

Equation (\ref{eq:BoyerSpectrum}) raises an immediate conceptual
question.

Why should a classical theory contain the quantum constant
$\hbar$?

SED answers this question differently from ordinary quantum mechanics.

Quantum mechanics assumes that $\hbar$ is a fundamental constant of
nature entering directly into the commutation relations,

\[
[x,p]=i\hbar.
\]

SED assumes instead that $\hbar$ merely determines the intensity of the
classical zero-point field.

From this perspective,

the appearance of Planck's constant is no more mysterious than the
appearance of the speed of light in Maxwell's equations.

Both are universal constants characterizing the vacuum.

%%%%%%%%%%%%%%%%%%%%%%%%%%%%%%%%%%%%%%%%%%%%%%%%%%%%%%%%%%%%%%%%%%%%%%%%%%%%%%%%
\subsection*{Critical Assessment}
%%%%%%%%%%%%%%%%%%%%%%%%%%%%%%%%%%%%%%%%%%%%%%%%%%%%%%%%%%%%%%%%%%%%%%%%%%%%%%%%

The mathematical construction of the zero-point field is internally
consistent and remarkably economical. Starting from classical Maxwell
theory, together with assumptions of homogeneity, isotropy and Lorentz
invariance, one is naturally led to a unique vacuum spectrum whose form
closely resembles that suggested by quantum theory.

Nevertheless, several important conceptual questions remain open.

The normalization of the spectrum cannot be derived within classical
electrodynamics alone and must be fixed empirically through the observed
value of Planck's constant.

Moreover, the total vacuum energy,

\[
\int_0^\infty
\frac{\hbar\omega^3}
{2\pi^2c^3}
\,d\omega,
\]

diverges quartically. In practice this divergence is handled by
considering energy differences, introducing physical cutoffs, or
restricting attention to observable correlation functions, much as in
quantum field theory. The physical interpretation of the absolute vacuum
energy remains one of the major unresolved issues shared by both SED and
conventional quantum electrodynamics.

From the perspective of the present monograph, the most important lesson
is that stochastic electrodynamics introduces randomness through the
vacuum field rather than through particle trajectories. This alternative
origin of stochasticity will provide an instructive point of comparison
with Nelson's diffusion theory and, later, with the stochastic
Hamiltonian formulation developed in the SHP framework.


%%%%%%%%%%%%%%%%%%%%%%%%%%%%%%%%%%%%%%%%%%%%%%%%%%%%%%%%%%%%%%%%%%%%%%%%%%%%%%%%
\section{Derivation of the Lorentz-Invariant Zero-Point Spectrum}
\label{sec:LorentzInvariantSpectrum}
%%%%%%%%%%%%%%%%%%%%%%%%%%%%%%%%%%%%%%%%%%%%%%%%%%%%%%%%%%%%%%%%%%%%%%%%%%%%%%%%

The cornerstone of stochastic electrodynamics is the assumption that the
vacuum contains a homogeneous, isotropic stochastic electromagnetic
field. Unlike ordinary thermal radiation, this field must exist even at
zero temperature and must not single out any preferred inertial frame.

These requirements are sufficiently restrictive that the spectral energy
density is determined almost uniquely. In this section we derive the
Lorentz-invariant spectrum following the physical arguments originally
developed by Boyer and later refined by several authors.

%%%%%%%%%%%%%%%%%%%%%%%%%%%%%%%%%%%%%%%%%%%%%%%%%%%%%%%%%%%%%%%%%%%%%%%%%%%%%%%%
\subsection{General Spectral Representation}
%%%%%%%%%%%%%%%%%%%%%%%%%%%%%%%%%%%%%%%%%%%%%%%%%%%%%%%%%%%%%%%%%%%%%%%%%%%%%%%%

Let

\[
u(\omega)\,d\omega
\]

denote the electromagnetic energy per unit volume contained in the
frequency interval

\[
(\omega,\omega+d\omega).
\]

The total vacuum energy density is therefore

\begin{equation}
u
=
\int_0^\infty
u(\omega)\,d\omega.
\end{equation}

At this stage the function

\[
\nu(\omega)
\]

is completely arbitrary.

Our task is to determine its form using only general physical
principles.

%%%%%%%%%%%%%%%%%%%%%%%%%%%%%%%%%%%%%%%%%%%%%%%%%%%%%%%%%%%%%%%%%%%%%%%%%%%%%%%%
\subsection{Dimensional Analysis}
%%%%%%%%%%%%%%%%%%%%%%%%%%%%%%%%%%%%%%%%%%%%%%%%%%%%%%%%%%%%%%%%%%%%%%%%%%%%%%%%

The dimensions of the spectral density are

\[
[\nu]
=
\frac{\mathrm{Energy}}
{\mathrm{Volume}\times\mathrm{Frequency}}.
\]

Using

\[
[\omega]=T^{-1},
\qquad
[c]=LT^{-1},
\]

one immediately finds that the only possible Lorentz-covariant
combination having the correct dimensions is

\begin{equation}
\nu(\omega)
=
C\,\omega^3,
\end{equation}

where

\[
C
\]

is a universal constant having dimensions of action divided by
\(c^3\).

Dimensional analysis alone therefore already suggests a cubic spectrum.

However, dimensional analysis cannot determine the constant
\(C\), nor can it prove Lorentz invariance.

A more detailed argument is required.

%%%%%%%%%%%%%%%%%%%%%%%%%%%%%%%%%%%%%%%%%%%%%%%%%%%%%%%%%%%%%%%%%%%%%%%%%%%%%%%%
\subsection{Lorentz Transformation of Plane Waves}
%%%%%%%%%%%%%%%%%%%%%%%%%%%%%%%%%%%%%%%%%%%%%%%%%%%%%%%%%%%%%%%%%%%%%%%%%%%%%%%%

Consider a monochromatic plane wave with four-wavevector

\[
k^\mu
=
\left(
\frac{\omega}{c},
\mathbf k
\right),
\qquad
k^\mu k_\mu=0.
\]

Under a Lorentz transformation,

\[
k^\mu
\longrightarrow
k'^\mu
=
\Lambda^\mu{}_\nu
k^\nu .
\]

The observed frequency changes according to the Doppler formula,

\begin{equation}
\omega'
=
\gamma
(1-\beta\cos\theta)\,
\omega .
\label{eq:Doppler}
\end{equation}

Consequently,

both the frequency and the density of states transform under boosts.

A Lorentz-invariant vacuum requires that these two effects compensate
each other exactly.

%%%%%%%%%%%%%%%%%%%%%%%%%%%%%%%%%%%%%%%%%%%%%%%%%%%%%%%%%%%%%%%%%%%%%%%%%%%%%%%%
\subsection{Density of Electromagnetic Modes}
%%%%%%%%%%%%%%%%%%%%%%%%%%%%%%%%%%%%%%%%%%%%%%%%%%%%%%%%%%%%%%%%%%%%%%%%%%%%%%%%

In a volume

\[
V,
\]

the number of electromagnetic modes contained in the interval

\[
(\omega,\omega+d\omega)
\]

is

\begin{equation}
g(\omega)d\omega
=
\frac{\omega^2}
{\pi^2c^3}
\,d\omega,
\label{eq:ModeDensity}
\end{equation}

including both polarization states.

The spectral energy density may therefore be written as

\begin{equation}
\nu(\omega)
=
g(\omega)\,
\varepsilon(\omega),
\end{equation}

where

\[
\varepsilon(\omega)
\]

is the average energy carried by each mode.

%%%%%%%%%%%%%%%%%%%%%%%%%%%%%%%%%%%%%%%%%%%%%%%%%%%%%%%%%%%%%%%%%%%%%%%%%%%%%%%%
\subsection{Scale Invariance of the Vacuum}
%%%%%%%%%%%%%%%%%%%%%%%%%%%%%%%%%%%%%%%%%%%%%%%%%%%%%%%%%%%%%%%%%%%%%%%%%%%%%%%%

At zero temperature the vacuum contains no intrinsic frequency scale.

The spectrum must therefore remain invariant under the simultaneous
transformation

\[
\omega
\rightarrow
\lambda\omega .
\]

Suppose

\[
\varepsilon(\omega)
\propto
\omega^n .
\]

Then

\[
\nu(\omega)
\propto
\omega^{n+2}.
\]

Lorentz invariance of the vacuum stress-energy tensor requires

\[
n=1,
\]

leading immediately to

\begin{equation}
\boxed{
\nu(\omega)
\propto
\omega^3.
}
\end{equation}

Thus the cubic spectrum is not an arbitrary assumption.

It follows directly from Lorentz symmetry together with the absence of
any preferred frequency scale.

%%%%%%%%%%%%%%%%%%%%%%%%%%%%%%%%%%%%%%%%%%%%%%%%%%%%%%%%%%%%%%%%%%%%%%%%%%%%%%%%
\subsection{Normalization}
%%%%%%%%%%%%%%%%%%%%%%%%%%%%%%%%%%%%%%%%%%%%%%%%%%%%%%%%%%%%%%%%%%%%%%%%%%%%%%%%

The proportionality constant cannot be determined by classical
electrodynamics.

Experiment fixes it to be

\begin{equation}
\boxed{
\nu(\omega)
=
\frac{\hbar\omega^3}
{2\pi^2c^3}.
}
\label{eq:FinalSpectrum}
\end{equation}

Equivalently,

the average energy per electromagnetic normal mode becomes

\begin{equation}
\boxed{
\varepsilon(\omega)
=
\frac12
\hbar\omega.
}
\label{eq:HalfQuantum}
\end{equation}

This expression is numerically identical to the zero-point energy of a
quantum harmonic oscillator.

In SED, however, its interpretation is entirely different.

It is regarded as the energy carried by each classical stochastic mode
of the electromagnetic vacuum.

%%%%%%%%%%%%%%%%%%%%%%%%%%%%%%%%%%%%%%%%%%%%%%%%%%%%%%%%%%%%%%%%%%%%%%%%%%%%%%%%
\subsection{Vacuum Correlation Functions}
%%%%%%%%%%%%%%%%%%%%%%%%%%%%%%%%%%%%%%%%%%%%%%%%%%%%%%%%%%%%%%%%%%%%%%%%%%%%%%%%

The spectral density determines all vacuum correlation functions.

For example,

\begin{equation}
\langle
E_i(\mathbf x,t)
E_j(\mathbf x',t')
\rangle
=
\int
d^3k\,
P_{ij}(\mathbf k)
\,
\nu(\omega)
\cos
\left[
\mathbf k\cdot
(\mathbf x-\mathbf x')
-
\omega(t-t')
\right],
\end{equation}

where

\[
P_{ij}
=
\delta_{ij}
-
\frac{k_ik_j}{k^2}
\]

is the transverse projection operator.

These correlation functions provide the stochastic forces responsible
for all particle dynamics in stochastic electrodynamics.

%%%%%%%%%%%%%%%%%%%%%%%%%%%%%%%%%%%%%%%%%%%%%%%%%%%%%%%%%%%%%%%%%%%%%%%%%%%%%%%%
\subsection{Comparison with Quantum Electrodynamics}
%%%%%%%%%%%%%%%%%%%%%%%%%%%%%%%%%%%%%%%%%%%%%%%%%%%%%%%%%%%%%%%%%%%%%%%%%%%%%%%%

The formal similarity between Eqs.~(\ref{eq:FinalSpectrum}) and
(\ref{eq:HalfQuantum}) and the corresponding expressions of quantum
electrodynamics is striking.

Nevertheless, the physical interpretations differ profoundly.

In quantum electrodynamics,

\[
\frac12\hbar\omega
\]

represents the ground-state energy of a quantized harmonic oscillator.

In stochastic electrodynamics,

the same quantity represents the average energy contained in a classical
random electromagnetic mode.

Thus identical mathematical expressions arise from two entirely
different ontological assumptions.

%%%%%%%%%%%%%%%%%%%%%%%%%%%%%%%%%%%%%%%%%%%%%%%%%%%%%%%%%%%%%%%%%%%%%%%%%%%%%%%%
\subsection*{Critical Assessment}
%%%%%%%%%%%%%%%%%%%%%%%%%%%%%%%%%%%%%%%%%%%%%%%%%%%%%%%%%%%%%%%%%%%%%%%%%%%%%%%%

The derivation above illustrates both the strength and the weakness of
stochastic electrodynamics.

On the one hand, the cubic spectrum follows naturally from very general
symmetry arguments and reproduces one of the most characteristic
numerical features of quantum vacuum fluctuations.

On the other hand, the overall normalization remains empirical. The
appearance of Planck's constant is not explained but assumed through the
measured intensity of the vacuum field. Consequently, SED shifts the
fundamental question from the quantization of matter to the origin of
the stochastic vacuum itself.

Nevertheless, the existence of a unique Lorentz-invariant classical
vacuum spectrum remains one of the most elegant and physically
compelling aspects of the theory.



%%%%%%%%%%%%%%%%%%%%%%%%%%%%%%%%%%%%%%%%%%%%%%%%%%%%%%%%%%%%%%%%%%%%%%%%%%%%%%%%
\section{Radiation Reaction and the Fluctuation--Dissipation Balance}
\label{sec:RadiationReaction}
%%%%%%%%%%%%%%%%%%%%%%%%%%%%%%%%%%%%%%%%%%%%%%%%%%%%%%%%%%%%%%%%%%%%%%%%%%%%%%%%

The existence of a stochastic electromagnetic vacuum is not, by itself,
sufficient to explain quantum phenomena. One must also understand how a
charged particle responds to this fluctuating field.

The dynamics of a charged particle differs fundamentally from ordinary
Brownian motion. Besides the random force exerted by the vacuum field,
an accelerated charge continuously emits electromagnetic radiation.
Energy is therefore lost through radiation reaction.

The motion is governed by the competition between two opposing effects:

\[
\boxed{
\text{Vacuum fluctuations}
\Longleftrightarrow
\text{Radiation reaction}.
}
\]

The former inject energy into the particle, whereas the latter removes
it. A stationary state becomes possible only if these two processes
balance each other.

%%%%%%%%%%%%%%%%%%%%%%%%%%%%%%%%%%%%%%%%%%%%%%%%%%%%%%%%%%%%%%%%%%%%%%%%%%%%%%%%
\subsection{Lorentz Force}
%%%%%%%%%%%%%%%%%%%%%%%%%%%%%%%%%%%%%%%%%%%%%%%%%%%%%%%%%%%%%%%%%%%%%%%%%%%%%%%%

The relativistic equation of motion is

\begin{equation}
m\frac{du^\mu}{d\tau}
=
eF^{\mu\nu}u_\nu ,
\label{eq:LorentzForce}
\end{equation}

where

\[
F^{\mu\nu}
=
F^{\mu\nu}_{\rm ext}
+
F^{\mu\nu}_{\rm ZP}.
\]

The second term denotes the stochastic zero-point field.

Equation (\ref{eq:LorentzForce}) is entirely classical.
Randomness enters only through the stochastic electromagnetic tensor.

%%%%%%%%%%%%%%%%%%%%%%%%%%%%%%%%%%%%%%%%%%%%%%%%%%%%%%%%%%%%%%%%%%%%%%%%%%%%%%%%
\subsection{Radiation Reaction}
%%%%%%%%%%%%%%%%%%%%%%%%%%%%%%%%%%%%%%%%%%%%%%%%%%%%%%%%%%%%%%%%%%%%%%%%%%%%%%%%

Accelerated charges radiate electromagnetic energy.

Conservation of energy therefore requires a self-force acting on the
particle.

In the nonrelativistic approximation,

\begin{equation}
m\ddot x
=
F_{\rm ext}
+
m\tau_0
\dddot x ,
\label{eq:AbrahamLorentz}
\end{equation}

where

\begin{equation}
\tau_0
=
\frac{2e^2}{3mc^3}.
\end{equation}

For an electron,

\[
\tau_0
\simeq
6.3\times10^{-24}\ {\rm s},
\]

an extremely short time scale.

The third-order derivative reflects the fact that radiation reaction
depends upon the rate of change of acceleration.

%%%%%%%%%%%%%%%%%%%%%%%%%%%%%%%%%%%%%%%%%%%%%%%%%%%%%%%%%%%%%%%%%%%%%%%%%%%%%%%%
\subsection{Relativistic Form}
%%%%%%%%%%%%%%%%%%%%%%%%%%%%%%%%%%%%%%%%%%%%%%%%%%%%%%%%%%%%%%%%%%%%%%%%%%%%%%%%

The fully covariant Abraham--Lorentz--Dirac equation is

\begin{equation}
m\frac{du^\mu}{d\tau}
=
eF^{\mu\nu}u_\nu
+
\frac{2e^2}{3c^3}
\left(
\frac{d^2u^\mu}{d\tau^2}
+
u^\mu
\frac{du_\nu}{d\tau}
\frac{du^\nu}{d\tau}
\right).
\label{eq:ALD}
\end{equation}

The second term is the radiation-reaction force.

Although derived entirely within classical electrodynamics, it has
generated considerable controversy because of its unusual mathematical
properties.

%%%%%%%%%%%%%%%%%%%%%%%%%%%%%%%%%%%%%%%%%%%%%%%%%%%%%%%%%%%%%%%%%%%%%%%%%%%%%%%%
\subsection{Runaway Solutions}
%%%%%%%%%%%%%%%%%%%%%%%%%%%%%%%%%%%%%%%%%%%%%%%%%%%%%%%%%%%%%%%%%%%%%%%%%%%%%%%%

Equation (\ref{eq:ALD}) possesses exponentially growing solutions,

\[
a(t)
\propto
e^{t/\tau_0},
\]

even when no external force acts.

These runaway solutions are clearly unphysical.

They arise because the equation is third order in time.

%%%%%%%%%%%%%%%%%%%%%%%%%%%%%%%%%%%%%%%%%%%%%%%%%%%%%%%%%%%%%%%%%%%%%%%%%%%%%%%%
\subsection{Preacceleration}
%%%%%%%%%%%%%%%%%%%%%%%%%%%%%%%%%%%%%%%%%%%%%%%%%%%%%%%%%%%%%%%%%%%%%%%%%%%%%%%%

A second difficulty is preacceleration.

Some solutions begin to accelerate before the external force is applied,

\[
a(t)\neq0,
\qquad
t<t_{\rm force}.
\]

This apparent violation of causality has stimulated an extensive
literature and remains one of the most discussed aspects of classical
electrodynamics.

%%%%%%%%%%%%%%%%%%%%%%%%%%%%%%%%%%%%%%%%%%%%%%%%%%%%%%%%%%%%%%%%%%%%%%%%%%%%%%%%
\subsection{Landau--Lifshitz Approximation}
%%%%%%%%%%%%%%%%%%%%%%%%%%%%%%%%%%%%%%%%%%%%%%%%%%%%%%%%%%%%%%%%%%%%%%%%%%%%%%%%

For slowly varying external fields one may eliminate these pathological
solutions perturbatively.

Replacing the acceleration appearing inside the radiation-reaction term
by the Lorentz force yields the Landau--Lifshitz equation,

\begin{equation}
m\frac{du^\mu}{d\tau}
=
eF^{\mu\nu}u_\nu
+
\frac{2e^3}{3mc^3}
\partial_\alpha F^{\mu\nu}
u^\alpha
u_\nu
+\cdots .
\end{equation}

This equation is second order and avoids both runaway solutions and
preacceleration within its domain of validity.

For most practical applications it has become the preferred classical
equation of motion.

%%%%%%%%%%%%%%%%%%%%%%%%%%%%%%%%%%%%%%%%%%%%%%%%%%%%%%%%%%%%%%%%%%%%%%%%%%%%%%%%
\subsection{Fluctuation--Dissipation Balance}
%%%%%%%%%%%%%%%%%%%%%%%%%%%%%%%%%%%%%%%%%%%%%%%%%%%%%%%%%%%%%%%%%%%%%%%%%%%%%%%%

The stochastic vacuum continuously transfers energy to the charged
particle,

\[
\left<
P_{\rm in}
\right>
>0.
\]

Radiation reaction removes energy,

\[
\left<
P_{\rm out}
\right>
>0.
\]

A stationary equilibrium requires

\begin{equation}
\boxed{
\left<
P_{\rm in}
\right>
=
\left<
P_{\rm out}
\right>.
}
\label{eq:FDT}
\end{equation}

This relation is the SED analogue of the fluctuation--dissipation
theorem familiar from statistical mechanics.

Unlike ordinary Brownian motion, however, both fluctuation and
dissipation originate from classical electrodynamics itself.

%%%%%%%%%%%%%%%%%%%%%%%%%%%%%%%%%%%%%%%%%%%%%%%%%%%%%%%%%%%%%%%%%%%%%%%%%%%%%%%%
\subsection{Relation to Brownian Motion}
%%%%%%%%%%%%%%%%%%%%%%%%%%%%%%%%%%%%%%%%%%%%%%%%%%%%%%%%%%%%%%%%%%%%%%%%%%%%%%%%

The analogy with Langevin dynamics is illuminating.

For an ordinary Brownian particle,

\[
m\dot v
=
-\gamma v
+\xi(t),
\]

where

\[
\xi(t)
\]

is the stochastic force and

\[
\gamma
\]

represents viscous damping.

In stochastic electrodynamics,

\[
\xi(t)
\longrightarrow
eE_{\rm ZP}(t),
\]

while the damping term is replaced by radiation reaction.

Thus the entire mathematical structure of Brownian motion reappears in a
purely electromagnetic context.

%%%%%%%%%%%%%%%%%%%%%%%%%%%%%%%%%%%%%%%%%%%%%%%%%%%%%%%%%%%%%%%%%%%%%%%%%%%%%%%%
\subsection{Conceptual Significance}
%%%%%%%%%%%%%%%%%%%%%%%%%%%%%%%%%%%%%%%%%%%%%%%%%%%%%%%%%%%%%%%%%%%%%%%%%%%%%%%%

The fluctuation--dissipation balance occupies a central position in
stochastic electrodynamics.

Without radiation reaction the stochastic vacuum would continuously
increase the particle energy.

Without vacuum fluctuations the particle would radiate away all of its
energy and eventually come to rest.

Only the simultaneous presence of both mechanisms permits stable
stationary states.

This balance is responsible for many of the characteristic predictions
of SED, including the ground-state energy of the harmonic oscillator.

%%%%%%%%%%%%%%%%%%%%%%%%%%%%%%%%%%%%%%%%%%%%%%%%%%%%%%%%%%%%%%%%%%%%%%%%%%%%%%%%
\subsection*{Critical Assessment}
%%%%%%%%%%%%%%%%%%%%%%%%%%%%%%%%%%%%%%%%%%%%%%%%%%%%%%%%%%%%%%%%%%%%%%%%%%%%%%%%

The interplay between vacuum fluctuations and radiation reaction is one
of the strongest aspects of stochastic electrodynamics because both
ingredients arise naturally within classical electrodynamics.

At the same time, radiation reaction remains one of the oldest
conceptual problems in theoretical physics. The Abraham--Lorentz--Dirac
equation exhibits runaway solutions and preacceleration, while the
Landau--Lifshitz equation is an approximation rather than an exact
theory.

Consequently, the foundations of SED inherit the unresolved difficulties
of classical radiation reaction. Whether these problems reflect a
limitation of classical electrodynamics itself or merely indicate the
need for a more complete microscopic description remains an open
question.

From the perspective of this monograph, an important lesson emerges:
the stochastic force and the dissipative self-force are not independent
ingredients but two complementary aspects of the interaction between a
charged particle and its electromagnetic environment. This observation
will later prove useful when constructing stochastic Hamiltonian
dynamics within the SHP framework.



